\subsection{1.1, Systems of linear equations}
\begin{problems}
  \problem Is the following equation linear in the variables $x_1$, $x_2$, and
  $x_3$?  Justify your answer.
  \[x_1 - 5 = x_2 - x_3.\]   %yes
  \begin{solution}
      Linear.  Can be written as $x_1 - x_2 + x_3 = 5$.
  \end{solution}

  \problem Is the following equation linear in the variables $x_1$, $x_2$, and
  $x_3$?  Justify your answer.
  \[x_1 - \sqrt[3]{x_2} = 1.\]
  \begin{solution}
    Not linear, due to the $\sqrt[3]{x_2}$ term.
  \end{solution}

  \problem Is the following equation linear in the variables $x_1$, $x_2$, and
  $x_3$?  Justify your answer.
    \[\ln(5) x_1 + x_2 - \sqrt{11} x_3 = 0.\]
    \begin{solution}
      Linear.  Here, $\ln(5)$ and $\sqrt{11}$ are real number coefficients of
      the $x_1$ and $x_3$ terms, respectively.
    \end{solution}
    
    \problem Is the following equation linear in the variables $x_1$, $x_2$, and
    $x_3$?  Justify your answer.
  \[e^{x_1} + e^{x_2} = e^{x_3}.\]
  \begin{solution}
      Not linear, due to variables being within exponents.
  \end{solution}

  \problem Is there a consistent system of three linear equations in two
  variables $x$ and $y$?  If your answer is yes, write such a system and draw a
  corresponding picture using $(x, y)$ to label points in ${\bf R}^2$.  If your
  answer is no, explain why no such system exists.  \newline
  \begin{solution}
    Yes, for example we could visualize three different lines in the
    $xy$-plane that have exactly one point in common.  
    If we want $(1, 1)$ to be that point, we could write equations for
    three different lines that contain $(1,1)$.
    \[x + 2y = 3\]
    \[x - y = 0\]
    \[2x + y = 3.\]
    \begin{center}
      \includegraphics[scale=0.5]{./1_1/lines.pdf}
    \end{center}      
  \end{solution}

  \problem In each case, determine whether the point is in the plane
  in ${\bf R}^3$ given by \[3x - 10y + z = 15.\] We follow the usual
  convention of using $(x, y, z)$ to label points in ${\bf R}^3$.
  \begin{subproblems}
    \subproblem $(3, -10, 1)$
    \subproblem $(3, -10, 1, 15)$
    \subproblem $(5, 1, 10)$
    \subproblem $(0, 0, 15)$
    \subproblem $(1, -1, 1)$
  \end{subproblems}
  \begin{solution}
    \begin{subproblems}
      \subproblem No: $3(3)-10(-10)+1 = 110$, not $15$. %
      \subproblem No: $(3, -10, 1, 15)$ is not even a point in
      ${\bf R}^3$. %
      \subproblem Yes: $3(5)-10(1)+10 = 15$. %
      \subproblem Yes: $3(0)-10(0)+15 = 15$. %
      \subproblem No: $3(1)-10(-1)+1 = 14$, not $15$.
    \end{subproblems}
  \end{solution}
  
  \problem Is there an inconsistent system of two linear equations in
  three variables $x$, $y$, and $z$?  If your answer is yes, write
  such a system and draw a corresponding picture using $(x, y, z)$ to
  label points in ${\bf R}^3$.  If your answer is no, explain why no
  such system exists. 
  \newline 
  \begin{solution}
    Yes, for example
    \[x + y + z = 0\]
    \[x + y + z = 10.\] %
    For this system to be consistent, we would need $0=10$, so the system
    is inconsistent. \newline \newline
    %We graph this below.
    %\begin{center}
    %  \includegraphics[scale=0.5]{./1_1/xyz.pdf}
    %\end{center}    
  \end{solution}

  \problem Consider the system of linear equations in $x$, $y$, and
  $z$ given by
  \[x + y + 0z = 0\]
  \[0x + 0y + z = 5.\]
  Describe its solution set geometrically, where we use $(x, y, z)$
  to label points in ${\bf R}^3$.
  \begin{solution}
    The solution set is a line.  It is the intersection of the
    horizontal plane $z=5$ with the vertical plane where $y= -x$.  
    %We
    %illustrate it graphically below.
    %\begin{center}
    %  \includegraphics[scale=0.4]{./1_1/intersection.png}
    %\end{center}
    We could parametrize the points on this line by $(x, -x, 5)$ where
    $x$ is real.
  \end{solution}
  

\problem Write a linear equation or a system of linear equations in
  the variables $x$, $y$, and $z$, so that the solution set is 
  the $xy$-plane. \newline
  \begin{solution}
      The $yz$-plane in space is given by the single linear equation 
      $x=0$, or in longhand: %
      \[x + 0y + 0z = 0.\] 
  \end{solution}

\problem Write a linear equation or a system of linear equations in
  the variables $x$, $y$, and $z$, so that the solution set is 
  a line. \newline
  \begin{solution}
      The $x$-axis in space the set of all points in space
      where $y=0$ and $z=0$, so our system is simply
      \[0x + y + 0z =0 \]
      \[0x + 0y + z = 0.\]
  \end{solution}

\problem Write a linear equation or a system of linear equations in
  the variables $x$, $y$, and $z$, so that the solution set is the 
  point $(1, 2, 3)$. \newline
  \begin{solution}
      The system of equations must only be satisfied when
      $x=1$, $y=2$, and $z=3$.  One way to do this is simply to make
      the three equations $x=1$, $y=2$, and $z=3$.  In long form,
      \[x + 0y + 0z = 1\]
      \[0x + y + 0z = 2\]
      \[0x + 0y + z = 3.\]      
  \end{solution}

\problem % Created for ILA
The following is a kind of common brain teaser that many people answer
incorrectly.  We will use linear algebra to solve it now. \newline
\newline A rare baseball card and a coffee mug cost $\$110$ together.
The baseball card costs $\$100$ more than the mug.  How much does the
baseball card cost?

\begin{solution}
  Let $x$ be the cost of the baseball card and $y$ be the cost of the
  mug.  We have been told that $x + y = 110$ and $x = y + 100$.  Now
  that we have written down this system of two linear equations in two
  variables, it is a standard high school algebra problem.  This gives
  $(y+100)+y= 110$ so $y=5$ and thus $x=105$.  Therefore, the baseball
  card costs $\$105$.
\end{solution}

\problem % Created for ILA
Suppose $f$ is a polynomial of the form $f(x) = ax^2+bx+c$
whose graph passes through the points $(-1, 11)$, $(2, 5)$, and
$(3, 7)$.  Write a system of linear equations that we could solve to
find $a$, $b$, and $c$.

\begin{solution}
  We use the facts $f(-1)=11$, $f(2) = 5$, and $f(3) = 7$:
  \[f(-1) = 11, \qquad a(-1)^2 + b(-1) + c = 11\]
  \[f(2) = 5, \qquad a(2^2) + b(2) + c = 5\]
  \[f(3) = 7 \qquad a(3^3) + b(3) + c = 7.\] This gives us the linear
  system of equations:
  \[a - b + c = 11\]
  \[4a + 2b + c = 5\]
  \[9a + 3b + c = 7.\]
\end{solution}

\problem % From AMS Notices, page 562, Volume 65, Number 5
This brain teaser comes from a memorial article for the mathematician
Paul Salley in the Notices of the American Mathematical Society,
Volume 65, Number 5.  The article is accessible at
{\color{blue}{\href{https://www.ams.org/journals/notices/201805/rnoti-p555.pdf?fbclid=IwAR1C1f6EtMygzK81ZGVzHwQFaLopa_O3sSQLXXDxYT-xvAldpmymLC75kEQ}{this
      link}}}.  %
This type of so-called ``Jim and John Problem'' has many variations
that sound like tricky brain teasers but can be solved by writing
systems of equations and using linear algebra.  We give a slight
variation here: \newline \newline Alice and Bob have the same
birthday.  Alice is 24 years old, which is twice as old as Bob was at
the time when Alice was as old as Bob is now.  How old is Bob now?

\begin{solution}
  While Alice and Bob are alive, the difference between their ages is
  some constant $c$ no matter what year we are in:
  \[(\textup{Alice's age}) - (\textup{Bob's age}) = c.\]
  We're told that alice is 24, so
  \[24-(\textup{Bob's age now}) = c.\]
  Also, from the information we are given, Bob was 12 when Alice was the same
  age that Bob is now.
  \[(\textup{Alice's age when Bob was 12}) - 12 = c,
    \quad \textup{ so} \quad (\textup{Bob's age now})-12 = c.\]
  We have a system of two equations: 
  \[24-(\textup{Bob's age now}) = c\]
  \[(\textup{Bob's age now}) - 12 = c.\]
  Subtracting the first equation from the second, we get
  \[-36 + 2(\textup{Bob's age now}) = 0,\] therefore Bob is now $18$ years
  old and $c = 6$ ($c=6$ means Alice is $6$ years older than Bob, which
  agrees with the fact that Alice is 24).
\end{solution}




%\problem % Created for ILA
%Iron, carbon, and oxygen can be combined to create steel in the
%chemical equation
%\[\textup{Fe}_2 \textup{O}_3 + 3 CO_2\]

\problem Consider a $3 \times 3$ QR code given by an array below,
where each spot can either have a 0 (for black) or 1 (for white).
\[\mat{\cdot \; \cdot \; \cdot;
    \cdot \; \cdot \; \cdot;
    \cdot \; \cdot \; \cdot}\]
\begin{subproblems}
  \subproblem Each QR code is an element in some ${\bf R}^n$.  What is
  $n$?
  %
  \subproblem How many $3 \times 3$ QR codes are there?
\end{subproblems}

\begin{solution}
  \begin{subproblems}
    \subproblem There are $3\times 3 = 9$ spots where we can put a $0$
    or a $1$, so $n = 9$.  \subproblem We have two choices for each
    entry in each QR code, with possibilities ranging from all zeros
    to all ones.  This gives us \[2^9 = 512\] such QR codes.
  \end{subproblems}
\end{solution}

\end{problems}\setcounter{problem}{0}



%
%
\newpage
%
%

\subsection{1.2, Row reduction}

\begin{problems}
\problem %Jankowski Spring 2018 exam
  True or false: The augmented matrix $\amat{0 1 0 2; 0 0 1 -3}$ is
  in reduced row echelon form. \newline
  \begin{solution}
      True.  The second and third columns are the pivot
      columns, and the matrix satisfies all conditions for RREF.
  \end{solution}
	
\problem %Jankowski Spring 2018 exam 
    True or false: If the reduced row echelon form of an augmented matrix has a row of 
    zeros, then the system of linear equations corresponding to the 
    augmented matrix must have infinitely many solutions.
    \newline
    \begin{solution}
        False.  For example, consider the augmented matrix
	\[\amat{0 0 1; 0 0 0}.\]
    \end{solution}

\problem Is the following matrix in reduced row echelon
  form (RREF)?  If not, row reduce to put it in RREF. %
  \[\amat{1 -2 4; 0 0 0}\]
  \begin{solution}
    The matrix is in RREF. %
  \end{solution}

\problem Is the following matrix in reduced row echelon
  form (RREF)?  If not, row reduce to put it in RREF. %
  \[\amat{0 1 2 3; 0 0 1 4}\]
  \begin{solution}
    The matrix is not in RREF because the pivot in the second line has a ``2'' in its column.  To put the matrix in RREF, we need to do one row operation: 
    \[\amat{0 1 2 3; 0 0 1 4}
    \xrightarrow[]{R_1 = R_1 - 2 R_2}
    \amat{0 1 0 -5; 0 0 1 4}.\]
  \end{solution}

\problem Is the following matrix in reduced row echelon
  form (RREF)?  If not, row reduce to put it in RREF. %
  \[\amat{1 0 0 1; 0 1 0 1; 0 0 1 0; 0 0 0 1}\]
  \begin{solution}
    No.  The rightmost column is a pivot column (with
      pivot in the fourth row), but the column has nonzero entries in
      its first and second row.
      \[\amat{1 0 0 1; 0 1 0 1; 0 0 1 0; 0 0 0 1}
        \xrightarrow[R_1 = R_1 - R_4]{R_2 = R_2 - R_4}
        \amat{1 0 0 0; 0 1 0 0; 0 0 1 0; 0 0 0 1}.\]  
  \end{solution}
  
  \problem Put the following matrix into reduced row
  echelon form (RREF).  Indicate which columns are pivot columns.
  $\mat{3 -9 0 3; 1 -3 1 -4; 4 -12 2 -6}$.

  \begin{solution}
    Try this problem using the Interactive Row Reducer
    \href{https://textbooks.math.gatech.edu/ila/demos/rrinter.html?mat=3,-9,0,3:1,-3,1,-4:4,-12,2,-6}{here}.

    Columns 1 and 3 are the pivot columns, and the RREF of the matrix
    is
    \[\mat{1 -3 0 1; 0 0 1 -5; 0 0 0 0}.\]
  \end{solution}

  \problem Put the following matrix into reduced row echelon form
  (RREF).  Indicate which columns are pivot columns.
  $\mat{-3 4 7; 2 1 -1; 3 1 -2; -2 2 4}$.
  
  \begin{solution}
    Try this problem using the Interactive Row Reducer
    \href{https://textbooks.math.gatech.edu/ila/demos/rrinter.html?mat=-3,4,7:2,1,-1:3,1,-2:-2,2,4}{here}.

    The quickest way might be to switch the first and fourth rows and
    then divide the new first row by $-2$ to make the numbers work
    nicely.
    
    Columns 1 and 2 are the pivot columns, and the RREF of the matrix
    is
    \[\mat{1 0 -1; 0 1 1; 0 0 0; 0 0 0}.\]
  \end{solution}

  \problem Put the following matrix $A$ into reduced row echelon form
  (RREF).  Indicate which columns are pivot columns.  Is the
  corresponding system of linear equations consistent or inconsistent?
  $\amat{5 10 -5 40; 4 -2 36 -28; 3 20 9 40}$.

  \begin{solution}
    Try this problem using the Interactive Row Reducer
    \href{https://textbooks.math.gatech.edu/ila/demos/rrinter.html?mat=5,10,-5,8:4,-2,36,-28:3,20,9,40}{here}.

    One way to do this problem is to divide the first row by $5$ and
    then do row operations to create zeros below the first entry of
    the left column.  Once we do this, we see that the second column
    is also pivot column.  Doing one more row operation gives us a row
    echelon form of the matrix, whereby we see all three columns to
    the left of the augment bar are pivot columns.  One row echelon form
    of the augmented matrix is
    \[\amat{1 2 -1 8; 0 1 -4 6; 0 0 68 -68}.\]

    We keep row-reducing to get the RREF of the augmented matrix:
    \[\amat{1 0 0 3; 0 1 0 2; 0 0 1 -1}.\]
    This system is consistent, with unique solution $x = 3$, $y = 2$,
    $z = -1$.
  \end{solution}
  


\problem Write an augmented matrix in RREF with three pivots that
  represents an inconsistent system of three equations in four
  variables.

	\begin{solution}
      Many possibilities.  For example,
      \[\amat{1 0 0 0 0; 0 1 0 0 0; 0 0 0 0 1},
        \qquad \amat{0 1 0 0 0; 0 0 1 0 0; 0 0 0 0 1}, \quad \textup{
          or } \quad \amat{0 0 1 0 0; 0 0 0 1 0; 0 0 0 0 1}.\]
	\end{solution}
	
  \problem Is there a system of three equations in two variables
    that has infinitely many solutions?  If your answer is yes, write
    an augmented matrix in RREF that represents one such system.

	\begin{solution}
      Yes.  For example, consider the system
	\[ x + y = 1\]
	\[ 2x + 2y = 2\]
	\[ 3x + 3y = 3.\] Each equation represents the line $x + y = 1$,
    so the solution set is all points $(x,y)$ on the line $y =
    1-x$. The corresponding augmented matrix and its RREF are
	\[\amat{1 1 1; 2 2 2; 3 3 3} \xrightarrow[]{RREF} \amat{1 1 1; 0 0
        0; 0 0 0}.\]
  \end{solution}

  \problem Row reduce the matrix to determine whether its
  corresponding system of linear equations is consistent.
  \[\amat{3 -2 11 2 5; 0 1 -1 2 3; 1 0 3 2 -1}.\]
  \begin{solution}
    We swap the first and third rows to begin.
    \begin{align*}
      \amat{3 -2 11 2 5; 0 1 -1 2 3; 1 0 3 2 -1}
      \xrightarrow[]{R_1 \leftrightarrow R_3}
      &
        \amat{1 0 3 2 -1; 0 1 -1 2 3; 3 -2 11 2 5}
      \\
      \xrightarrow[]{R_3 = R_3 - 3 R_1}
      &
        \amat{1 0 3 2 -1; 0 1 -1 2 3; 0 -2 2 -4 8}
      \\
      \xrightarrow[]{R_3 = R_3 + 2R_2}
      &
        \amat{1 0 3 2 -1; 0 1 -1 2 3; 0 0 0 0 14}
    \end{align*}
    We do not need to row-reduce any further because we have found
    a pivot in the rightmost column, therefore the system is inconsistent.
    The final row of the augmented matrix corresponds to the equations
    ``$0 = 14$.''
  \end{solution}
  
  \problem Are there three lines in the $xy$-plane (not necessarily
  distinct lines) whose common intersection is exactly one point?  If
  so, give an example and write an augmented matrix for the
  corresponding system of linear equations in $x$ and $y$.
	
  \begin{solution}
    Yes.  For example, take the lines $x - y = 0$, $2x-2y = 0$, and
    $y =1$.  The first two lines are the same line, which intersects
    the non-parallel line $y=1$ at exactly one point, namely when
    $x=y=1$. To illustrate the example, we solve the augmented system
    below to find the unique solution.
	\begin{align*}
      \amat{1 -1 0; 2 -2 0;  0 1 1} 
      & \xrightarrow[]{R_2 = R_2 -2R_1}
        \amat{1 -1 0; 0 0 0; 0 1 1} \xrightarrow[]{R_1 \leftrightarrow R_2}
        \amat{1 -1 0; 0 1 1; 0 0 0} \xrightarrow[]{R_1 = R_1 + R_2} 
        \amat{1 0 1; 0 1 1; 0 0 0}.
	\end{align*}
  \end{solution}

  \problem Write the following system of equations in augmented matrix
  form, and put the matrix into RREF.  After this, write two solutions
  $(x_1, x_2, x_3)$ 
  to the system of equations.
  \[2x_1 - 3x_2 + x_3 = 6\]
  \[x_1 - x_2 = 4\]
  \[x_2 - x_3 = 2.\]
  \begin{solution}
    We put this into an augmented matrix and row-reduce:
    \begin{align*}
      \amat{2 -3 1 6; 1 -1 0 4; 0 1 -1 2}
      &
        \xrightarrow[]{R_1 \leftrightarrow R_2}
        \amat{1 -1 0 4; 2 -3 1 6; 0 1 -1 2}
        \xrightarrow[]{R_2 = R_2 - 2 R_1}
        \amat{1 -1 0 4; 0 -1 1 2; 0 1 -1 2}
      \\
      &
        \xrightarrow[\textup{then } R_3 = R_3 - R_2]{R_2 = (-1)*R_2}
        \amat{1 -1 0 4; 0 1 -1 2; 0 0 0 0}
        \xrightarrow[]{R_1 = R_1 + R_2}
        \amat{1 0 -1 6; 0 1 -1 2; 0 0 0 0}.
    \end{align*}
    Now that we have found the RREF of the matrix, we see that
    $(x_1, x_2, x_3)$ is a solution as long as $x_1 - x_3 = 6$ and
    $x_2 - x_3 = 2$.  To find a solution, we could just pre-set our
    favorite value for one of the variables and see if we could solve
    for the others.  For example, if we set $x_1 = 1$, then
    to satisfy $x_1 - x_3 = 6$ we need $x_3 = -5$, and now the second
    equation becomes
    \[x_2 - x_3 = 2, \qquad x_2 - (-5) = 2, \qquad x_2 = -3.\]
    Therefore, $(1, -3, -5)$ is a solution.

    To get a second solution, we could set $x_3 = 0$ and solve for $x_1$
    and $x_2$:  $x_1 - x_3 = 6$ gives $x_1 = 6$, and then $x_2-x_3 = 2$
    gives $x_2 = 2$, so $(6, 2, 0)$ is a second solution.

    In reality, there are infinitely many solutions to this system of
    equations.  We'll see in section 1.3 that this consistent system
    has $x_3$ as a \emph{free} variable because its column does not
    have a pivot, which means that we can set $x_3$ equal to any real
    number and then solve for $x_1$ and $x_2$.
  \end{solution}

  \problem Find all values of $h$ so that the following system has at least one solution. % Dan's notes
  \[x + y = 1\]
  \[2x + 2y = h.\]
  \begin{solution}
      The answer is $h=2$ only.  Subtracting two times the first equation from the second equation gives us
      \[x+y = 1\]
      \[0 = h - 2.\]
      The system will be consistent as long as $h-2 = 0$ (so $h=2$).  In this case, the system has infinitely many solutions since both equations will describe the same line $x+y = 1$. 
      
      If $h \neq 2$, then the equation $0= h-2$ will not be satisfied, so the system is inconsistent.
  \end{solution}

  \problem Extra practice with row reduction that you can check using the interactive row reducer.  Put the following matrix into reduced row echelon form (RREF). % Dan's notes
  \[\amat{1 2 3 9; 2 -1 1 8; 3 0 -1 3}\]
  \begin{solution}
      You can check your work using the interactive row reducer 
      {\color{blue} \href{https://textbooks.math.gatech.edu/ila/demos/rrinter.html?mat=1,2,3,9:2,-1,1,8:3,0,-1,3}{here}}.
      The RREF is
      \[\amat{1 0 0 2; 0 1 0 -1; 0 0 1 3}.\]
  \end{solution}

  \problem Extra practice with row reduction that you can check using the interactive row reducer.  Put the following matrix into reduced row echelon form (RREF). % Dan's notes
  \[\amat{0 7 -4 2; 2 4 6 12; 3 1 -1 -2}\]
  \begin{solution}
      You can check your work using the interactive row reducer 
      {\color{blue} \href{https://textbooks.math.gatech.edu/ila/demos/rrinter.html?mat=0,7,-4,2:2,4,6,12:3,1,-1,-2}{here}}.
      The RREF is
      \[\amat{1 0 0 -5/9; 0 1 0 10/9; 0 0 1 13/9}.\]
  \end{solution}

 \problem Extra practice with row reduction that you can check using the interactive row reducer.  Put the following matrix into reduced row echelon form (RREF). % Dan's notes
  \[\amat{4 -5 3 2; 1 -1 -2 -6; 4 -4 -14 18}\]
  \begin{solution}
      You can check your work using the interactive row reducer 
      {\color{blue} \href{https://textbooks.math.gatech.edu/ila/demos/rrinter.html?mat=4,-5,3,2:1,-1,-2,-6:4,-4,-14,18}{here}}.
      The RREF is
      \[\amat{1 0 0 -123; 0 1 0 -103; 0 0 1 -7}.\]
  \end{solution}
\end{problems} \setcounter{problem}{0}



\newpage

\subsection{1.3, Parametric form}

\begin{problems}
  \problem % Dan's notes
  True or false: If a system of linear equations has 100 variables and 70
  equations, then there must be infinitely many solutions.
  
\begin{solution}
  False.  It might be inconsistent, for example one equation could be $0 = 1$.
  It might also have infinitely many solutions.  For example, if all 70
  equations are each just the equation
  \[x_1 + x_2 + x_3 + \cdots + x_{100} = 0.\] The corresponding augmented matrix
  will not have a pivot in its rightmost column, and it will have 99 non-pivot
  columns to the left of the augment bar.
\end{solution}


\problem % Dan's notes
True or false: If a system of linear equations has 70 variables and 100
equations, then it must be inconsistent.

\begin{solution}
  False.  It might be consistent, for example a system where all 100 equations
  are just repeats of the equation
    \[x_1 + x_2 + \cdots + x_{70} = 0.\]
\end{solution}



\problem % Dan's notes
\begin{subproblems}
  \subproblem How can we tell if an augmented matrix corresponds to a consistent
  system of linear equations?  \subproblem Suppose a consistent system of linear
  equations has finitely many solutions.  How many solutions must it have?
\end{subproblems}
\begin{solution}
\begin{subproblems}
  \subproblem The system is consistent if and only if the rightmost column of
  the augmented matrix is NOT a pivot column.  %
  \subproblem It must have exactly one solution.  Since the system is
  consistent, there is no pivot in the rightmost column of the augmented matrix,
  leaving us with two possibilities: Either the augmented matrix has a pivot in
  every column to the left of the augment bar (no free variables, thus exactly
  one solution), or there is a non-pivot column to the left of the augment bar
  (free variable, thus infinitely many solutions).
\end{subproblems}
\end{solution}

\problem   %Jankowski Spring 2018 exam 1 
Is it possible for a system of three linear equations in four
variables to have exactly one solution?  Justify your answer.

\begin{solution}
  No.  This system will correspond to an augmented matrix that has 3
  rows and 4 columns to the left of the augment.  This means the
  matrix has at most 3 pivots, so at least one variable will be free
  and the system will have infinitely many solutions if it is
  consistent.
\end{solution}


\problem Write the augmented matrix corresponding to the system of
equations and put the matrix in reduced row echelon form.  Write the
set of solutions in parametric form.  Indicate which variables are
free variables.
  \[-5x_1 + 6x_2 - 5x_3 = 17\]
  \[2x_2 = 4\]
  \[-2x_1 + 2x_2 - 2x_3 = 6\]

\begin{solution}
  \begin{align*}
    \amat{-5 6 -5 17; 0 2 0 4; -2 2 -2 6}
    & 
      \xrightarrow[]{R_1 \leftrightarrow R_3}
      \amat{-2 2 -2 6; 0 2 0 4; -5 6 -5 17}
      \xrightarrow[]{R_1 = \frac{-R_1}{2}}
      \amat{1 -1 1 -3; 0 2 0 4; -5 6 -5 17}
    \\
    &
      \xrightarrow[]{R_3= R_3 + 5R_1} \amat{1 -1 1 -3; 0 2 0 4; 0 1 0 2}
      \xrightarrow[]{R_2 = \frac{R_2}{2}} \amat{1 -1 1 -3; 0 1 0 2; 0 1 0 2}
    \\
    &
      \xrightarrow[]{R_3 = R_3 - R_2}
      \amat{1 0 1 -1; 0 1 0 2; 0 0 0 0} 
  \end{align*}
  The first two columns are the pivot columns, so the system is
  consistent and $x_3$ is free.  The first equation is
  $x_1 + x_3 = -1$ and the second is $x_2 = 2$, so
  \[x_1 = -1-x_3, \qquad x_2 = 2, \qquad x_3 = x_3 \quad \textup{
      ($x_3$ any real number)}.\] Equivalently, \
  $(x_1, x_2, x_3) = (-1 - x_3, \ 2, \ x_3) \ \textup{ where
    $x_3$ is any real number}.$
\end{solution}

\problem % Created for ILA
Write the augmented matrix corresponding to the system of
equations and put the matrix in reduced row echelon form.  Write the
set of solutions in parametric form.  Indicate which variables are
free variables.
\[x_1 - 2x_2 -x_3 + x_5 = 3\]
  \[x_3 - 2x_4 + 3x_5 = 1\]
\begin{solution}
\[
  \amat{1 -2 -1 0 1 3; 0 0 1 -2 3 1} \xrightarrow[]{R_1 = R_1+R_2}
  \amat{1 -2 0 -2 4 3; 0 0 1 -2 3 1}.\]
%
The first and third columns are the pivot columns, so the system is
consistent and the variables $x_2$, $x_4$, and $x_5$ are free.
Putting the free variables to the right hand side in the equations
$x_1 - 2x_2 - 2x_4 + 4x_5 = 3$ and $x_3-2x_4+3x_5 = 1$, we get the
equations below where $x_2$, $x_4$, and $x_5$ are any real numbers:
\[
  x_1 = 3 + 2x_2 + 2x_4 -4x_5 \qquad x_2 = x_2 \qquad x_3 = 1 + 2x_4 -
  3x_5 \qquad x_4 = x_4 \qquad x_5 = x_5.
\]
Equivalently,
\[(x_1, x_2, x_3, x_4, x_5) = (3 + 2x_2 + 2x_4 -4x_5, \ x_2, \ 1 +
  2x_4 - 3x_5, \ x_4, \ x_5)\] where $x_2$, $x_4$, and $x_5$ are any
real numbers.
\end{solution}

\problem % created for ILA
Write the augmented matrix corresponding to the system of
equations and put the matrix in reduced row echelon form.  Write the
set of solutions in parametric form.  Indicate which variables are
free variables.
\[x_1 + x_2 - 2x_3 + 3x_4 = 4\]
\[-2x_1 + 3x_2 + x_4 = 1\]
\[3x_1 + 8x_2 - 10x_3 + 16x_4 = 21\]
\[-5x_1 + 5x_2 + 2x_3 - x_4 = -2.\]
\begin{solution}
  \begin{align*}
    \amat{1 1 -2 3 4; -2 3 0 1 1; 3 8 -10 16 21; -5 5 2 -1 -2}
    \xrightarrow[R_3 = R_3 -3R_1, \ R_4 = R_4 + 5R_1]{R_2 = R_2 + 2R_1}
    &
      \amat{1 1 -2 3 4; 0 5 -4 7 9; 0 5 -4 7 9; 0 10 -8 14 18}
    \\
    \xrightarrow[R_4 = R_4 - 2R_2]{R_3 = R_3 - R_2}
    &                                                  
      \amat{1 1 -2 3 4; 0 5 -4 7 9; 0 0 0 0 0; 0 0 0 0 0}
    \\
    \xrightarrow[]{R_2 = R_2/5}
    &
      \amat{1 1 -2 3 4; 0 1 -4/5 7/5 9/5; 0 0 0 0 0; 0 0 0 0 0}
    \\
    \xrightarrow[]{R_1 = R_1 - R_2}
    &
      \amat{1 0 -6/5 8/5 11/5; 0 1 -4/5 7/5 9/5; 0 0 0 0 0; 0 0 0 0 0}.
  \end{align*}
  The system is consistent, with free variables $x_3$ and $x_4$:
  \[x_1 = \dfrac{11}{5} + \dfrac{6}{5} x_3 - \dfrac{8}{5} x_4
    \qquad
    x_2 = \dfrac{9}{5} + \dfrac{4}{5} x_2 - \dfrac{7}{5} x_4
    \qquad
    x_3 = x_3
    \qquad
    x_4 = x_4.
    \]
\end{solution}

\problem We have seen that any system of linear equations with more
than one solution must have infinitely many solutions.  Give an
example of a system of two equations in two variables that is not
linear and has exactly two solutions.

\begin{solution}
  Picture curves that are described naturally by fundamental equations
  in precalculus and intersect at two points.  For example, the curves
  $y = \sqrt{x}$ and $y=x^2$ intersect at exactly two points, namely
  $(0, 0)$ and $(1,1)$.  One way to write this system is
	\[\sqrt{x} - y = 0\]
	\[x^2 - y = 0.\]	
\end{solution}

\problem Is it possible for a system of four linear equations in four
unknowns to have infinitely many solutions?  If your answer is yes,
write such a system.  If your answer is no, justify why it is not
possible.
	
\begin{solution}
  Yes.  We could simply write the same linear equation four times, but
  this is not the only possibility.  For example,
  \[x_1 + x_2 + x_3 = 0\]
  \[2x_1 + 2x_2 + x_3 = 0\]
  \[2x_1 + 2x_2 + x_3 = 0\]
  \[x_3 = 1\]
\end{solution}

\problem Bruce Hernandez is infatuated with the system of linear
equations given by
  \[3x - 2y = 4\]
  \[6x + h y = k,\] where $h$ and $k$ are some real numbers.

\begin{subproblems}
  \subproblem Determine all values of $h$ and $k$ (if there are any)
  so that the system of equations is inconsistent.%
  \subproblem Determine all values of $h$ and $k$ (if there are any)
  so that the system of equations has infinitely many solutions.
\end{subproblems}

\begin{solution}
  \[ \amat{3 -2 4; 6 h k} \xrightarrow[]{R_2=R_2-2R_1} \amat{3 -2 4; 0
      h+4 k-8}.\]
\begin{subproblems}
  \subproblem The system is inconsistent if and only if the right
  column is a pivot column.  Therefore, we need $h+4 = 0$ and
  $k-8 \neq 0$ so that the rightmost entry of the second row is a
  pivot.
  \[ h= -4 \qquad k \neq 8.\] %
  \subproblem To have infinitely many solutions, we must have a free
  variable AND the rightmost column cannot be a pivot colum, so the
  entire second row above must be zero.  Thus $h+4=0$ and $k-8=0$.
  \[h= -4 \qquad k = 8.\] %
  *** In both parts, we are finding $h$ and $k$ so that the lines
  given by each equation are parallel.  In (a), we solve so that we
  get two \textit{different} parallel lines.  In (b), each equation
  gives the \textit{same} line since the second row is two times the
  first row.
\end{subproblems}
\end{solution}

\problem % created for ILA
Write a system of linear equations in the variables
$x_1, x_2, x_3, x_4$ whose general solution has parametric form
\[x_1 = 5 + x_3 \qquad
  x_2 = 10 - x_3 \qquad
  x_3 = x_3 \qquad
  x_4 = -7.
\]
\begin{solution}
  We have exactly 1 free variable in a consistent linear system in 4
  variables, so we must have 3 pivots in our corresponding augmented
  matrix.  This means we must have at least 3 rows in our augmented
  matrix.  One possibility is
  \[\amat{1 0 -1 0 5; 0 1 1 0 10; 0 0 0 1 -7}.\]  
\end{solution}

\problem Find every polynomial $f$ of degree $2$ whose graph passes
through the points $(-1, 11)$, $(2, 5)$, and $(3, 7)$.

\begin{solution}
  With $f(x) = ax^2 + bx + c$ (where $a \neq 0$ since $f$ has degree
  $2$), and using the facts $f(-1)=11$, $f(2) = 5$, and $f(3) = 7$:
  \[a(-1)^2 + b(-1) + c = 11\]
  \[a(2^2) + b(2) + c = 5\]
  \[a(3^3) + b(3) + c = 7.\] This gives us a linear system of
  equations, which we put into an augmented matrix and row-reduce:
  \[a - b + c = 11\]
  \[4a + 2b + c = 5\]
  \[9a + 3b + c = 7.\]
  %
  \[\amat{1 -1 1 11; 4 2 1 5; 9 3 1 7}
    \xrightarrow[]{RREF}
    \amat{1 0 0 1; 0 1 0 -3; 0 0 1 7}.\]
  Therefore $a=1$, $b= -3$, and $c=7$, so there is exactly one
  degree $2$ polynomial $f$ that goes through the three points given in the problem:
  \[f(x) = x^2 - 3x + 7.\]
\end{solution}

\problem Find every polynomial $g$ of degree $3$ whose graph passes
through the points $(0, 1)$, $(1, 1)$, and $(2, 9)$.
%Compare your answer to the previous problem.

\begin{solution}
  We let $g(x) = ax^3 + bx^2 + cx + d$ (where $a \neq 0$) and plug in
  the three points:
  \[a(0^3)+b(0^2)+c(0)+d = 1\]
  \[a(1^3)+b(1^2)+c(1)+d = 1\]
  \[a(2^3)+b(2^2)+c(2)+d = 9.\] We put this into a system and
  row-reduce its corresponding augmented matrix:
  %
  \[d = 1\]
  \[a + b + c + d = 1\]
  \[8a + 4b + 2c + d = 9.\]
  %
  \[\amat{0 0 0 1 1; 1 1 1 1 1; 8 4 2 1 9}
    \xrightarrow[]{RREF} \amat{1 0 -1/2 0 2; 0 1 3/2 0 -2; 0 0 0 1
      1}.\] Thus $c$ is free and
  \[a = 2 + \frac{c}{2} \quad b = -2 - \frac{3c}{2}, \quad d = 1.\]

  This gives us infinitely degree $3$ polynomials $g$ going through
  the three prescribed points, depending upon which value we choose
  for $c$ (we just need to choose $c \neq -4$ so that $a \neq 0$,
  otherwise $g$ will have degree less than $3$).  For example, when
  $c = 1$ we get the polynomial
  \[g_1(x) = \frac{3}{2} x^3 - \frac{7}{2}x^2 + x + 1\]
  whereas when $c = 2$ we get
  \[g_2(x) = 3x^3 - 5x^2 + 2x + 1.\]
  We could even take $c = 0$ and get
  \[g_3(x) = 2x^3 - 2x^2 + 1,\]
  and so on.
\end{solution}

\end{problems}\setcounter{problem}{0}

%
%
\newpage
%
%

